\section{Paper2Agent 评测与 Agent-to-Agent 协作}

有了流程设计，还要看证据。本节聚焦两类结果：一是 Paper2Agent 相对直接``把论文+repo 喂给通用 coding agent''的性能差异；二是多个 Paper Agent 协作时是否能生成新见解。

\subsection{性能与稳健性：先建稳固 MCP，再做下游推理}

课程中给出的结论是：先构建 robust paper MCP，再执行任务，整体准确性与稳定性优于直接让通用 coding agent 现场摸索 repo。并且在展示案例中，Paper2Agent runtime 约为对照方式的三分之一以内，说明预处理阶段的结构化投资能够在推理阶段回收。

\begin{importantbox}{为什么``预处理重''反而可能``总时长更短''}
\begin{itemize}
    \item 直接执行时错误多、回滚多、重试多，累计成本高。
    \item MCP 先把环境和接口固定，减少推理时不必要分支。
    \item 复用同一 MCP 时，后续任务边际成本快速下降。
\end{itemize}
\end{importantbox}

\subsection{Agent-to-Agent 协作：方法 Agent 与数据 Agent 自动对接}

Zou 描述了一个有代表性的场景：某团队发布方法论文，另一团队发布数据论文。过去需要作者间大量沟通才能完成对接；现在可让两个 Paper Agent 基于各自 MCP 先自动探索``方法是否适配该数据''，再把高价值候选提交给人类复核。

课程举例里，AlphaGenome 相关方法 Agent 与 ADHD GWAS 数据 Agent 协作后，识别到新的潜在线索（splicing error 与风险关联）。这个例子最值得注意的并非结论本身，而是机制：\textbf{跨论文协作从社会网络驱动，变为协议驱动}。

\begin{knowledgebox}{Agent-to-Agent 科研协作的四步模板}
\begin{itemize}
    \item 能力声明：每个 Agent 公开可做什么、不能做什么。
    \item 任务对齐：定义共同目标与评价指标。
    \item 协作执行：交换中间结果并触发互相调用。
    \item 人类复核：对高影响结论做统计与实验双重验证。
\end{itemize}
\end{knowledgebox}

\subsection{概念漂移与边界控制}

问答环节有人追问：底层模型有预训练知识，如何避免超出论文证据边界的``概念漂移''？课程回答非常清楚：关键不在于禁止外部知识，而在于用文档完备、可测试、可追责的 paper MCP 作为主执行边界，并在不满足条件时触发降级或拒答。

\begin{warningbox}{如果没有边界控制，Paper Agent 会变成``会说但不保真''}
\begin{itemize}
    \item 对外看起来流畅，实际引用与实现脱节。
    \item 对代码细节的错误想当然会污染后续分析链。
    \item 在多 Agent 协作中错误会跨 Agent 传播并放大。
\end{itemize}
\end{warningbox}

\subsection{本章小结}

Paper2Agent 的实用价值建立在两点：稳健 MCP 带来的工程收益，以及多 Agent 协作带来的知识组合收益。前者提升效率与可靠性，后者提升科研搜索空间。

